problem:  
Let \((M,\Delta,g)\) be an *equiregular contact* sub-Riemannian manifold endowed with its Hausdorff measure \(\mu\).  
Assume that at every point \(x\in M\), the metric-measure tangent cone \(\operatorname{Tan}(M,d,\mu,x)\) is isomorphic to a fixed Carnot group \((G,d_G,\mu_G)\) with Haar measure, and that \((M,d,\mu)\) satisfies a sub-Riemannian measure-contraction property \(MCP(K,N)\) (in the sense of Juillet, Rizzi, and Agrachev).  

For such contact structures it is known (Agrachev–Barilari–Rizzi) that the small-ball volume admits an asymptotic expansion  
\[
V(x,r):=\mu(B(x,r)) = C(x)\,r^{Q}\Bigl(1 + \kappa(x)\,r^2 + o(r^2)\Bigr),
\]
where \(Q\) is the Hausdorff dimension, \(C(x)\) depends on the nilpotent approximation, and \(\kappa(x)\) encodes the horizontal curvature invariants of \((M,\Delta,g)\).  

**Open problem:**  

1. Within this class of contact, equiregular sub-Riemannian manifolds satisfying \(MCP(K,N)\), is the coefficient \(\kappa(x)\) uniquely determined (up to normalization) by the *metric-measure structure* \((M,d,\mu)\) alone?  
Equivalently, can two contact sub-Riemannian manifolds \((M_1,\Delta_1,g_1)\) and \((M_2,\Delta_2,g_2)\) have the same metric \(d\), the same Hausdorff measure \(\mu\), and the same tangent Carnot group (so identical \(C(x)\)) but differ in curvature coefficient \(\kappa(x)\)?  

2. If such deformation freedom exists, what is the smallest metric-measure invariant beyond the order \(r^{Q+2}\) expansion capable of detecting differences in horizontal curvature (for example, the \(r^{Q+4}\) or heat-kernel coefficients)?  

Why is it a "good" problem:  
This formulation explicitly restricts to contact sub-Riemannian manifolds where the small-ball expansion is known to exist, removing technical ambiguity while keeping the geometric question profound. It asks whether *second-order curvature information*—already known to appear in analytic expansions—can be *reconstructed purely from metric-measure data*.  

Conceptually, the problem probes the boundary between **microlocal curvature analysis** and **synthetic metric-measure geometry**, asking whether curvature is intrinsically encoded in the volume growth of infinitesimal metric balls, as it is in Riemannian geometry.  

A positive answer would demonstrate a powerful rigidity: the metric-measure data of an equiregular contact space fully determine its curvature invariants, potentially leading to a synthetic definition of sub-Riemannian curvature analogous to Ricci curvature in the \(MCP\) setting.  
A negative answer would reveal the existence of nontrivial deformations invisible to all first- and second-order metric-measure data, forcing the introduction of new invariants for sub-Riemannian spaces.  

The problem is *simple to state* (about the coefficient of an asymptotic expansion), *deep in implications*, and *bridges multiple fields*—synthetic geometry, sub-Riemannian analysis, and geometric measure theory—making it a quintessential "narrow entrance, profound depth" research challenge.